\section{Identification: contagion, common shocks and under-reporting}\label{sec:identification}

Incidents cluster in time. The substantive question is \emph{why}. Clustering may
reflect a salient external event whose effect lasts for days (a shock with an extended
response). It may reflect incidents triggering further incidents (contagion). Or it
may reflect both. Inference about contagion must also cope with two features of
incident data: many shocks are never recorded in any event calendar, and many incidents
are never recorded at all (the dark figure of Paper~1). This section gives four sets of
results.
\begin{enumerate}[(1)]
  \item Without a shock calendar, and for exponential or completely monotone kernels
    with no excite-only streams, the autocorrelation of counts cannot distinguish
    contagion from clustered exogenous shocks (\cref{thm:equiv,prop:th-cm}). Delayed
    kernels and excite-only streams do leave second-order traces
    (\cref{rem:th-noncm,prop:th-mv}). Omitting shocks creates apparent contagion
    (\cref{thm:spurious}).
  \item Identification is restored by a recorded shock calendar, by higher-order
    information, or by provenance links (\cref{cor:equiv-id}).
  \item Independent under-recording attenuates apparent contagion in closed form, which
    permits a dark-figure correction (\cref{thm:thinning,cor:th-darkfig,prop:th-hetthin}).
  \item Online-to-offline Granger non-causality is equivalent to a zero cross-excitation
    kernel, with caveats (\cref{prop:granger}).
\end{enumerate}
The tension between ``true'' and ``apparent'' contagion is old. \Citet{feller1943}
showed that the negative binomial law arises both from a contagion urn and from a
Poisson law with a heterogeneous rate. \Cref{thm:equiv} is a dynamic, second-order
analogue for self-exciting processes. Full proofs are in \cref{app:idproofs}.

\paragraph{Conventions.} $H(\mu,n,\beta)$ denotes the stationary univariate Hawkes
process with kernel $\varphi(t)=n\beta e^{-\beta t}$, $n\in(0,1)$, and mean
$\bar\Lambda=\mu/(1-n)$ (\cref{prop:stationary}). Two stationary point processes are
\emph{second-order equivalent} if they have the same mean and the same Bartlett
spectrum. $\Pr_X$ denotes the law of a process $X$, and $F_\infty=2\pi f(0)/\bar\Lambda$
is the asymptotic index of dispersion (\cref{sec:th-second}).

\subsection{Second-order observational equivalence}\label{sec:th-equiv}

Unrecorded shocks are modelled as random.

\begin{definition}[Shot-noise Cox process]\label{def:th-sncox}
Let $\nu\ge0$ and $r,a,\gamma>0$. Let $\{c_\ell\}$ be a homogeneous Poisson process
on $\mathbb R$ with rate $r$; these are the times of unrecorded shocks. Define
\[
\xi(t)=\nu+\sum_{\ell:\,c_\ell\le t}a\,e^{-\gamma(t-c_\ell)} .
\]
Conditionally on $\xi$, let $N$ be a Poisson process with intensity $\xi$. We write
$N\sim C(\nu,r,a,\gamma)$ \citep{cox1955,moller2003}.
\end{definition}

The process $C$ has no contagion: given the shocks, incidents occur independently.
Each shock produces $\Pois(m)$ incidents, $m=a/\gamma$, at independent
$\mathrm{Exp}(\gamma)$ delays, so $C$ is also a Neyman--Scott cluster process with mean
$\nu+r m$. Any stationary Cox process directed by $\xi$ has Bartlett spectrum
$(2\pi)^{-1}[\bar\Lambda+\hat c_\xi(\omega)]$, where $c_\xi$ is the covariance function
of $\xi$ and $\hat c_\xi\ge0$ (\cref{lem:th-cox2}). By Campbell's theorem,
$c_\xi(u)=\{r a^2/(2\gamma)\}e^{-\gamma|u|}$ for $C$, so
\begin{equation}\label{eq:th-fC}
  f_C(\omega)=\frac1{2\pi}\Big[\nu+\frac{r a}{\gamma}+\frac{r a^2}{\gamma^2+\omega^2}\Big].
\end{equation}

\begin{theorem}[Second-order observational equivalence]\label{thm:equiv}
Let $H=H(\mu,n,\beta)$ and $C=C(\nu,r,a,\gamma)$.
\begin{enumerate}[(a)]
  \item $H$ and $C$ are second-order equivalent if and only if
  \begin{equation}\label{eq:th-match}
    \gamma=\beta(1-n),\qquad \nu+\frac{r a}{\gamma}=\bar\Lambda,\qquad
    r a^2=\bar\Lambda\,\beta^2 n(2-n).
  \end{equation}
  \item For given $H$, the solutions of \eqref{eq:th-match} form a one-parameter
  family, indexed by the shock amplitude $a\in[a_{\min},\infty)$ with
  $a_{\min}=\beta n(2-n)/(1-n)$:
  \[
    \gamma=\beta(1-n),\qquad
    r=\frac{\bar\Lambda\beta^2 n(2-n)}{a^2},\qquad
    \nu=\bar\Lambda\Big(1-\frac{a_{\min}}{a}\Big).
  \]
  Equivalently, the mean number of incidents per shock, $m=a/\gamma$, satisfies
  $m\ge m_{\min}=n(2-n)/(1-n)^2=F_\infty-1$, and $r m^2=\bar\Lambda(F_\infty-1)$.
  The lower bound $a=a_{\min}$ gives $\nu=0$: every incident is then shock-driven.
  \item For any matched pair and any bin width $\Delta>0$, the binned counts of $H$
  and of $C$ have the same mean and autocovariance function. Hence they also share
  their variance-time curves, index of dispersion and expected periodograms. As
  $T\to\infty$, every continuous function of finitely many empirical means and
  autocovariances of binned counts has the same probability limit under $H$ and
  under $C$.
\end{enumerate}
\end{theorem}

The extra degree of freedom $a$ is invisible at second order. It trades off frequent,
small shocks ($a\downarrow a_{\min}$) against rare, large ones ($a\to\infty$, with
$\nu\to\bar\Lambda$). The shock response decays at rate $\beta(1-n)$, the inverse
of the \emph{cluster} time scale of $H$, rather than at the kernel rate $\beta$. Thus a
Hawkes
fit with branching ratio $n$ is matched, in its autocorrelations, by a world without
contagion in which unrecorded shocks each trigger on average at least
$(1-n)^{-2}-1$ incidents. The full laws of $H$ and $C$ do differ
(\cref{cor:equiv-id}(ii)).

\paragraph{General kernels.} The exponential kernel is not special. Completely monotone
kernels have the form $g(t)=\int_{(0,\infty)}e^{-bt}\sigma(\dd b)$ for a positive
measure $\sigma$. They include finite mixtures of exponentials and the Lomax (Omori-type) kernels
$g(t)=(\vartheta/\varsigma)(1+t/\varsigma)^{-\vartheta-1}$ with
$\vartheta,\varsigma>0$.

\begin{proposition}[Completely monotone kernels]\label{prop:th-cm}
Let $H$ be the stationary Hawkes process with background $\mu$ and kernel
$\varphi=ng$, where $n\in(0,1)$ and $g$ is a completely monotone probability density
on $(0,\infty)$ with $g(0+)<\infty$.
\begin{enumerate}[(a)]
  \item The resolvent $R=\sum_{j\ge1}\varphi^{*j}$ is completely monotone. The
  covariance density of $H$ has the form $c_H(u)=\int e^{-b|u|}\,\mathcal K(\dd b)$ for a
  positive measure $\mathcal K$ with $Z:=\int \mathcal K(\dd b)/b=\bar\Lambda(F_\infty-1)/2$. If $g$ is a
  mixture of $K$ exponentials, then $\mathcal K$ has at most $K$ atoms.
  \item Fix $m\ge F_\infty-1=n(2-n)/(1-n)^2$. Let $\{c_\ell\}$ be Poisson with rate
  $r=2Z/m^2$, and give the shocks independent decay rates $\Gamma_\ell$ with law
  $\pi(\dd b)=\mathcal K(\dd b)/(bZ)$. The Cox process directed by
  \[
  \xi(t)=\nu+\sum_{\ell:\,c_\ell\le t}m\,\Gamma_\ell\,e^{-\Gamma_\ell(t-c_\ell)},
  \qquad \nu=\bar\Lambda-r m\ \ (\ge0),
  \]
  is second-order equivalent to $H$. Each shock produces $\Pois(m)$ incidents with an
  exponential response at a random rate, and there is no contagion.
  \item In particular, $f_H(\omega)\ge\bar\Lambda/(2\pi)$ for every $\omega$.
\end{enumerate}
\end{proposition}

\begin{remark}[Delayed excitation leaves a second-order trace]\label{rem:th-noncm}
By \cref{lem:th-cox2}, every Cox process satisfies $f\ge\bar\Lambda/(2\pi)$. Since
$2\pi f_H(\omega)/\bar\Lambda=|1-\hat\varphi(\omega)|^{-2}$, a Hawkes process with
$2\operatorname{Re}\hat\varphi(\omega)<|\hat\varphi(\omega)|^2$ at some $\omega$ is
second-order distinguishable from \emph{every} Cox process. This happens whenever
$\operatorname{Re}\hat g$ becomes negative, as it does for delayed or hump-shaped
kernels. An example is the gamma kernel $g(t)\propto t^{k-1}e^{-\beta t}$ with $k>1$,
for which $\hat g(\omega)=(1+i\omega/\beta)^{-k}$. For $k=2$ a direct calculation gives
\[
\min_\omega 2\pi f_H(\omega)/\bar\Lambda=1-n/4,
\]
attained at $\omega=\beta\sqrt{3-n}$. The dip lies at frequencies of the order of the
inverse delay, so it is easily lost when counts are binned coarsely. For kernels that
are not completely monotone but still satisfy $f_H\ge\bar\Lambda/(2\pi)$, whether a
second-order-equivalent Cox process always exists remains open.
\end{remark}

\paragraph{Several streams.} If $\Phik$ is diagonal, the streams are independent, and
\cref{thm:equiv,prop:th-cm} apply to each stream with its own shock process. Cross-excitation
is different. Write $\bm E(\omega)=2\pi\bm f(\omega)-\operatorname{diag}(\bar{\bm\Lambda})$
for the \emph{excess spectral matrix}.

\begin{proposition}[Multivariate streams]\label{prop:th-mv}
\begin{enumerate}[(a)]
  \item For every stationary $M$-variate Cox process with square-integrable driver and
  integrable covariance, $\bm E(\omega)=\hat{\bm c}_\xi(\omega)$ is Hermitian and
  positive semidefinite for every $\omega$.
  \item For the Hawkes process of \cref{prop:stationary},
  $\bm E_H(\omega)=\bm G(\omega)\operatorname{diag}(\bar{\bm\Lambda})\bm G(\omega)^{*}-\operatorname{diag}(\bar{\bm\Lambda})$
  with $\bm G=(\Imat-\hat\Phik)^{-1}$. Suppose some stream $k$ receives no excitation
  ($\varphi_{kl}\equiv0$ for all $l$) but excites another stream $j$
  ($\alpha_{jk}>0$). Then $\bm E_H(\omega)$ is not positive semidefinite for $\omega$
  near $0$, and no $M$-variate Cox process, with any common or idiosyncratic shocks, is
  second-order equivalent to it.
  \item Let $\{c_\ell\}$ be a homogeneous Poisson process of rate $r$ (common
  unrecorded shocks), with i.i.d.\ random loadings $\bm a_\ell\in[0,\infty)^M$,
  independent of $\{c_\ell\}$ and with $\E|\bm a_\ell|^2<\infty$. Let the response
  shapes $s_k\ge0$ be square-integrable and stream-specific, and let $\bm a$ denote a
  generic loading:
  \[
  \xi_k(t)=\nu_k+\sum_\ell a_{\ell k}\,s_k(t-c_\ell).
  \]
  The cross-covariance is
  $c_{kj}(u)=r\,\E[a_ka_j]\int s_k(t+u)s_j(t)\,\dd t$, where $\bm a=(a_k)$. It is symmetric in $u$ when
  $s_k=s_j$, but asymmetric in general. Lead--lag cross-correlation between two streams
  is therefore not, by itself, evidence of directional excitation.
\end{enumerate}
\end{proposition}

\subsection{Spurious contagion from unrecorded shocks}\label{sec:th-spurious}

Now reverse the direction. Suppose that the truth has no contagion and that the
analyst fits a Hawkes process with a constant background.

\begin{theorem}[Spurious contagion]\label{thm:spurious}
Let $N$ be a stationary Cox process directed by a stationary $\xi\ge0$ with mean
$\bar\Lambda>0$, $\E\xi(0)^2<\infty$ and integrable covariance $c_\xi\not\equiv0$.
Define the \emph{burstiness index}
\[
\kappa=\frac1{\bar\Lambda}\int_{\mathbb R}c_\xi(u)\,\dd u=\frac{\hat c_\xi(0)}{\bar\Lambda}\ \ge0,
\qquad\text{so that }F_\infty=1+\kappa .
\]
\begin{enumerate}[(a)]
  \item \emph{Exact projection.} If $c_\xi(u)=\sigma^2e^{-\gamma|u|}$, then exactly one
  exponential Hawkes process is second-order equivalent to $N$. Its parameters are
  \begin{equation}\label{eq:th-nstar}
    n^{*}=1-(1+\kappa)^{-1/2},\qquad \beta^{*}=\gamma(1+\kappa)^{1/2},\qquad
    \mu^{*}=\bar\Lambda(1+\kappa)^{-1/2},
  \end{equation}
  where $\kappa=2\sigma^2/(\bar\Lambda\gamma)$. For $C(\nu,r,a,\gamma)$,
  $\kappa=r a^2/(\bar\Lambda\gamma^2)$. The second-order structure of a two-state
  regime-switching rate is also exponential: if $\xi$ switches between $\lambda_0$
  and $\lambda_1$ at rates $q_{01}$ and $q_{10}$, with
  $\pi_1=q_{01}/(q_{01}+q_{10})$, then
  $\kappa=2(\lambda_1-\lambda_0)^2\pi_0\pi_1/\{\bar\Lambda(q_{01}+q_{10})\}$.
  \item \emph{Moment-based estimators.} In the setting of (a), let $N$ be ergodic. Let
  $\hat\theta_T=G(\hat{\bm m}_T)$, where $\hat{\bm m}_T$ collects the empirical mean and
  finitely many empirical autocovariances of the binned counts on $[0,T]$, and $G$ is
  continuous. If $\hat\theta_T$ is consistent for $(\mu,n,\beta)$ under every
  stationary exponential Hawkes process, then $\hat\theta_T\to(\mu^*,n^*,\beta^*)$ in
  probability under $N$.
  \item \emph{Any covariance shape.} The dispersion functional
  $n_F=1-F_\infty^{-1/2}$ equals $n$ for every stationary Hawkes process, whatever its
  kernel. Applied to $N$, it equals $n^*=1-(1+\kappa)^{-1/2}$, which is positive
  whenever $\kappa>0$.
  \item \emph{Likelihood.} Now let $N$ be any stationary point process with mean
  $\bar\Lambda$, $\E N(0,1]^2<\infty$ and covariance density $c$; for a Cox process,
  $c=c_\xi$. Fit the family
  $\lambda_\theta(t)=\mu+n\int_{(-\infty,t)}g_\beta(t-s)\,\dN(s)$ with
  $\theta=(\mu,n,\beta)\in(0,\infty)\times[0,1)\times\mathcal B$, where each $g_\beta$
  is a probability density on $(0,\infty)$ with $\int_0^\infty g_\beta|c|<\infty$. Let
  \[
  L(\theta)=\E\int_0^1\log\lambda_\theta(t)\,\dN(t)-\E\lambda_\theta(0)
  \]
  be the expected log-likelihood per unit time. If $\int_0^\infty g_{\beta_0}(u)c(u)\,\dd u>0$
  for some $\beta_0\in\mathcal B$, then
  \[
  \sup_{\theta:\,n>0}L(\theta)>\sup_{\theta:\,n=0}L(\theta)=\bar\Lambda\log\bar\Lambda-\bar\Lambda .
  \]
  Hence every maximiser of $L$, that is, every pseudo-true value of the misspecified
  maximum likelihood estimator, has $n>0$. For $C(\nu,r,a,\gamma)$ with exponential
  $g_\beta$, $\int g_\beta c=r a^2\beta/\{2\gamma(\beta+\gamma)\}>0$ for every $\beta$.
\end{enumerate}
\end{theorem}

\begin{remark}[Reading \cref{thm:spurious}]\label{rem:th-spurious}
\begin{enumerate}[(i)]
  \item \emph{Interpreting $\kappa$.} For $C$, $\kappa=s\,m$, where $s=r m/\bar\Lambda$
  is the share of incidents triggered by unrecorded shocks and $m$ is the mean number
  of incidents per shock. Hence $n^*=1-(1+sm)^{-1/2}$. As a purely arithmetical
  illustration, $s=0.3$ and $m=10$ give $n^*=1/2$. Rare, large shocks and slow regime
  switches give large $\kappa$, and so apparent near-criticality ($n^*\uparrow1$). This
  is consistent with the ``apparent criticality'' that \citet{filimonov2015} discuss
  for Hawkes fits to data with non-stationary activity.
  \item \emph{What is and is not claimed.} Parts (a)--(c) concern second-order
  estimators, such as variance-time and autocovariance fits. Spectral (Whittle-type)
  estimators from binned counts \citep{cheysson2022} use the whole periodogram, so they
  are not covered by (b) as stated. Their population contrast
  $\int\{\log f_\theta+f/f_\theta\}$ is minimised exactly where $f_\theta=f$, which in
  case (a) happens only at $\theta^*$. Under the usual M-estimation conditions, which we
  do not verify, they therefore also converge to $\theta^*$. Part (d) proves
  only that the likelihood pseudo-true value has $n>0$. That value depends on the law
  of $N$ beyond second order, so it need not equal $n^*$. Nothing forces the members of
  the matched family of \cref{thm:equiv}, which all share $n^*$, to share it. The
  simulation study (\cref{sec:simulation}) compares maximum likelihood fits with $n^*$.
  Under the usual conditions for misspecified likelihoods \citep{white1982,ogata1978},
  the estimator converges to the maximisers of $L$; we do not verify those conditions
  here.
  \item \emph{Consequence.} Salient events that are missing from the calendar, or whose
  response is mis-specified, are absorbed into $\Amat$. Contagion estimates should
  therefore be reported together with the shock calendar, and with their sensitivity to
  omitting events from it.
\end{enumerate}
\end{remark}

\subsection{Restoring identification}\label{sec:th-restore}

\begin{corollary}[Identification restored]\label{cor:equiv-id}
\begin{enumerate}[(i)]
  \item \emph{Recorded shocks.} Consider model \eqref{eq:th-intensity} on $[0,T]$,
  started from the empty history, under \cref{ass:th-model}, with the calendar
  $(\tau_e,S_e)$ observed. Let
  $m_k(t)=\mu_k(t)+\sum_e\beta_kS_e\psi_k(t-\tau_e)$ be the immigrant intensity. The
  law of $\bm N$ on $[0,T]$ determines $m_k$ almost everywhere on $[0,T]$, and each
  $\varphi_{kj}$ almost everywhere on $(0,T)$. Two further conditions complete
  identification of all parameters:
  \begin{itemize}
    \item each kernel family is determined by its restriction to $(0,T)$, as
      exponential, gamma and power-law families are;
    \item the map $(\theta_k,\beta_k,\psi_k)\mapsto m_k$ is injective. The data on
      $[\tau_1,\tau_2)$ identify only the product $\beta_kS_1\psi_k$ there. Injectivity
      therefore holds, for instance, when $\tau_1>0$, $S_1>0$, $\mu_k$ is constant
      (or analytic in $t$ with linearly independent basis functions), and the
      \emph{scaled} family $\{c\,\psi:c>0,\psi\in\Psi_k\}$ is determined by its
      restriction to an interval. An example is $\Psi_k$ consisting of functions that
      are analytic on $(0,\infty)$; then $\int\psi_k=1$ pins down $\beta_k$. A
      family such as $\psi_s=s^{-1}\mathbf 1_{[0,s]}$ fails this condition.
  \end{itemize}
  Spline backgrounds are not analytic, so this sufficient condition does not cover the
  working model of \cref{ass:th-model}(i). There, separating background from shock
  response rests on the shock responses not being representable in the spline space,
  for example through a jump at $\tau_e$ when $\psi_k(0+)>0$ and no knot sits at
  $\tau_e$. We do not formalise this.
  In particular, when the $g_{kj}$ do not vanish on $(0,T)$, $\Amat=\zero$ is
  distinguishable from every $\Amat\ne\zero$: clustering explained by recorded shocks
  is not confounded with contagion.
  \item \emph{Higher-order information.} Let $C=C(\nu,r,a,\gamma)$ be arbitrary.
  Then $\Pr_H\neq\Pr_C$. Counts over long windows already reveal this. The limits
  $\chi_j=\lim_{T\to\infty}T^{-1}\operatorname{cum}_j N(0,T]$ are
  $\mu\E\mathcal S^j$ ($\mathcal S\sim$ Borel$(n)$) for $H$ and $\nu+r\E[\Pois(m)^j]$ for $C$.
  Matching $\chi_1,\chi_2,\chi_3$ determines $(\nu,r,m)$
  uniquely, with $\gamma$ free. With $\gamma=\beta(1-n)$, this choice is a member of
  the family of \cref{thm:equiv}. For every such $C$,
  \[
  \chi_4^{H}-\chi_4^{C}=\frac{\mu\,n^3\,Q(1-n)}{(1-n)^7(2-n)}>0,\qquad
  Q(u)=2u^4+8u^3+15u^2+16u+6 ,
  \]
  so no shot-noise Cox process matches the first four cumulant rates of $H$.
  \item \emph{Provenance.} Suppose each true parent--offspring pair is recorded as a
  link independently with probability $q\in(0,1]$, and there are no false links. Links
  might come, for example, from perpetrator-confirmed imitation, a shared perpetrator
  or a shared campaign. In the stationary multivariate model, the expected number of
  links per unit time from stream $j$ to stream $k$ is
  $q\,\alpha_{kj}\bar\Lambda_j$. Hence the support of $\Amat$ (the contagion graph) is
  identified even when $q$ is unknown, and $\Amat$ is identified when $q$ is known.
\end{enumerate}
\end{corollary}

Part (i) is the basis of the shock-augmented likelihood in \cref{sec:estimation}. Once
shocks are recorded, the exogenous explanation becomes the nested sub-model
$\Amat=\zero$, and the boundary likelihood-ratio test of $\alpha_{kj}=0$ applies. The
result is relative to the calendar: unrecorded shocks remain subject to
\cref{thm:spurious}. Hawkes models with exogenous news events in the intensity are
used in finance for the same reason \citep{rambaldi2015}. Part (ii) is a
statement of principle. Fourth-order cumulants are expensive to estimate, and the
practical vehicle for higher-order information is the full likelihood.

\begin{remark}[Why intensity jumps do not settle the question]\label{rem:th-filter}
It is sometimes said that a Hawkes intensity jumps at events whereas a Cox intensity
does not. This is true only of the Cox intensity \emph{given} the latent driver $\xi$,
which is unobservable. With respect to the observed history, the Cox intensity is the
filter $\lambda^N(t)=\E[\xi(t)\mid\Hist^N_{t-}]$. Wherever $\xi$ is continuous, Bayes'
rule shows that $\lambda^N$ jumps at each event by
$\Var(\xi(t)\mid\Hist^N_{t-})/\E(\xi(t)\mid\Hist^N_{t-})\ge0$ \citep[cf.][]{bremaud1981}.
Learning about an unobserved shock looks like self-excitation.

What does separate the models is the \emph{form} of the filter. For a Hawkes process
the observable intensity is affine in the past. \Cref{lem:th-affine} shows that, among
all stationary laws second-order equivalent to $H$, $\Pr_H$ is the only one whose
observable intensity is affine in the past. Here ``affine'' means
$\lambda(t)=c+\int_{(-\infty,t)}h(t-s)\,\dN(s)$ with $c>0$, $h\ge0$,
$h\in L^1\cap L^2$ and $\int h<1$. By (ii),
the filter of every matched $C$ is therefore nonlinear. The information that separates
contagion from shocks lies in this nonlinearity.
\end{remark}

\subsection{Under-reporting masks contagion}\label{sec:th-thinning}

Paper~1 models recorded incidents as an independent thinning of all incidents, with
recording probability $p$ equal to the reciprocal of the dark-figure ratio.

\begin{theorem}[Under-reporting masks contagion]\label{thm:thinning}
Let $H=H(\mu,n,\beta)$, and let $N^p$ be its independent $p$-thinning: each incident
is recorded independently, with probability $p\in(0,1]$, independently of $H$.
\begin{enumerate}[(a)]
  \item $N^p$ is stationary with mean $p\bar\Lambda$ and Bartlett spectrum
  \[
  f_p(\omega)=\frac{p\bar\Lambda}{2\pi}\Big[1+\frac{p\,\beta^2 n(2-n)}{\beta^2(1-n)^2+\omega^2}\Big].
  \]
  \item $N^p$ is second-order equivalent to exactly one exponential Hawkes process,
  $H(\mu',n',\beta')$, with
  \[
  \beta'(1-n')=\beta(1-n),\qquad \beta'^2=\beta^2(1-n)^2+p\,\beta^2n(2-n),\qquad
  \mu'=p\bar\Lambda(1-n'),
  \]
  \begin{equation}\label{eq:th-hp}
    n'=h_p(n)=1-\frac{1-n}{\sqrt{(1-n)^2+p\,n(2-n)}} .
  \end{equation}
  Equivalently, $F'_\infty-1=p\,(F_\infty-1)$: thinning scales the excess dispersion by
  $p$. We write $n'$ for the observed (apparent) branching ratio, called $n_{\mathrm{obs}}$
  in the overview of contributions.
  \item For $p\in(0,1)$, $h_p$ is a strictly increasing bijection of $[0,1)$, and
  $pn<h_p(n)<n$ for $n\in(0,1)$. Moreover $h_p(n)=pn+O(n^2)$ as $n\downarrow0$,
  $h_p(n)\to1$ as $n\uparrow1$, $h_p(n)\to0$ as $p\downarrow0$, $h_1(n)=n$, and
  $\beta(1-n)\le\beta'\le\beta$. The inverse is
  \begin{equation}\label{eq:th-hinv}
    h_p^{-1}(n')=1-\Big[1+\frac{(1-n')^{-2}-1}{p}\Big]^{-1/2}.
  \end{equation}
  \item For $p<1$, $N^p$ is not a Hawkes process, for any kernel and any constant
  background. Thus $n'$ is an \emph{apparent} branching ratio in a precise sense. It is
  the branching ratio of the unique exponential Hawkes process that is second-order
  equivalent to $N^p$. It is therefore the probability limit under $N^p$ of every
  moment-based estimator as in \cref{thm:spurious}(b). We do not identify the
  pseudo-true value of the Hawkes maximum likelihood estimator under $N^p$, which
  need not equal $n'$.
\end{enumerate}
\end{theorem}

Attenuation is never worse than the naive factor $p$, and it is weakest near
criticality. Recording reduces each individual offspring count by the factor $p$, but
the dispersion of whole clusters, which is what second-order statistics see, is
reduced less.

\begin{corollary}[Dark-figure correction]\label{cor:th-darkfig}
Let $\hat n'$ estimate $n'$ from recorded incidents. Let $\hat p$ estimate the
recording probability from Paper~1 ($p=1/D$, where $D$ is the dark-figure ratio), and
let $\hat p$ be independent of $\hat n'$. Set $\hat n=h_{\hat p}^{-1}(\hat n')$ using
\eqref{eq:th-hinv}. If $\hat n'\to n'$ and $\hat p\to p$, then $\hat n\to n$. By the
delta method,
\begin{equation}\label{eq:th-delta}
  \Var(\hat n)\approx\Big[\frac{(1-n)^3}{p\,(1-n')^3}\Big]^2\Var(\hat n')
  +\Big[\frac{n(1-n)(2-n)}{2p}\Big]^2\Var(\hat p),
  \qquad \Var(\hat p)\approx\frac{\Var(\hat D)}{D^4}.
\end{equation}
\end{corollary}

The second term in \eqref{eq:th-delta} does not shrink as $T$ grows: uncertainty about
coverage puts a floor under the precision of any dark-figure-corrected contagion
estimate. A Bayesian alternative passes posterior draws of $p$ from Paper~1 through
$h_p^{-1}$. If $\hat n'$ is the Hawkes maximum likelihood estimator, its limit is the
pseudo-true value under $N^p$, so the correction is exact only to the extent that this
value coincides with $n'$ (\cref{thm:thinning}(d)). The simulation study quantifies the
gap.

\paragraph{Heterogeneous recording.} Severe incidents may be both more likely to be
recorded and more likely to trigger further incidents. Suppose each incident carries an
\emph{unpredictable} mark $X\in\{1,\dots,K\}$, such as a severity class. Marks are
i.i.d.\ with law $\pi$, independent of the past. An incident with mark $x$ has expected
direct offspring $n_x$ and is recorded independently with probability $p_x$, and the
kernel is $\beta e^{-\beta t}$. Write $\E$, $\Var$ and $\Cov$ for moments under $\pi$,
and set $\bar n=\E n_X<1$ and $\bar p=\E p_X$.

\begin{proposition}[Heterogeneous recording]\label{prop:th-hetthin}
The recorded process is stationary with mean $\bar p\bar\Lambda$, where
$\bar\Lambda=\mu/(1-\bar n)$, and Bartlett spectrum
\[
f_{\mathrm{rec}}(\omega)=\frac{\bar p\bar\Lambda}{2\pi}\Big[1+\frac{\beta^2\{2(1-\bar n)\E[p_Xn_X]+\bar p\,\E[n_X^2]\}}{\beta^2(1-\bar n)^2+\omega^2}\Big].
\]
This is the spectrum of an exponential Hawkes process with branching ratio
\[
n''=1-\frac{1-\bar n}{\sqrt{(1-\bar n)^2+2(1-\bar n)\E[p_Xn_X]+\bar p\,\E[n_X^2]}} .
\]
The homogeneous correction of \cref{cor:th-darkfig} with $p=\bar p$ returns
$n_{\mathrm{naive}}=h_{\bar p}^{-1}(n'')$, where
\[
(1-n_{\mathrm{naive}})^{-2}=(1-\bar n)^{-2}\Big[1+\frac{2(1-\bar n)\Cov(p_X,n_X)}{\bar p}+\Var(n_X)\Big].
\]
It overstates $\bar n$ whenever $2(1-\bar n)\Cov(p_X,n_X)/\bar p+\Var(n_X)>0$, in
particular when more productive incidents are recorded more often.
\end{proposition}

With $p_x\equiv p$ and $n_x\equiv n$, \cref{prop:th-hetthin} reduces to
\cref{thm:thinning}. Heterogeneous productivity alone ($p\equiv1$) already inflates the
apparent branching ratio of unmarked data, to $1-(1-\bar n)/\sqrt{1+\Var(n_X)}$. The
remedy is to model marks explicitly, with mark-specific productivities, and to use
category-specific coverage when it is available. With stream-specific recording
probabilities $\bm P=\operatorname{diag}(p_k)$, the excess spectral matrix becomes
$\bm P\bm E_H\bm P$. This need not be of multivariate exponential-Hawkes form, and
the univariate correction cannot simply be applied stream by stream, because cascades
pass through unrecorded incidents. The multivariate case calls for an explicit thinning
model.

\subsection{Granger causality between online and offline streams}\label{sec:th-granger}

Let $V=[M]$. Stream $j$ does \emph{not Granger-cause} stream $k$, relative to $V$, if
the intensity $\lambda_k$ of $N_k$ with respect to $\Hist^V$ has a version that is
predictable with respect to $\Hist^{V\setminus\{j\}}$, the history generated by the
shock calendar and the other streams. This is local independence in the sense of
\citet{didelez2008}, the continuous-time analogue of \citet{granger1969}.

\begin{proposition}[Granger non-causality; \citealp{eichler2017}]\label{prop:granger}
Under \eqref{eq:th-intensity} and \cref{ass:th-model}, with the process started from
the empty history on $[0,T]$, stream $j$ does not Granger-cause stream $k$ if and only
if $\varphi_{kj}=0$ almost everywhere on $(0,T)$. For exponential kernels, this holds
if and only if $\alpha_{kj}=0$.
\end{proposition}

\Citet{eichler2017} prove the corresponding equivalence for stationary Hawkes processes
and build graphical models on it. \Cref{app:idproofs} gives a short proof for the
shock-augmented model. With $k=\mathrm{off}$ and $j=\mathrm{on}$, the hypothesis
$\alpha_{\mathrm{off},\mathrm{on}}=0$ states that online volume carries no predictive
information about offline incidents beyond that in the other modelled streams and the
calendar. This is the hypothesis tested by the boundary-corrected likelihood-ratio test
of \cref{sec:estimation}.

\begin{remark}[Caveats for online-to-offline inference]\label{rem:th-granger-caveats}
\begin{enumerate}[(i)]
  \item \emph{Common shocks.} Granger causality is relative to the information set.
  A salient event that is missing from the calendar, and that affects online volume
  faster than offline incidents, induces online-to-offline Granger causality in the
  modelled system even if $\alpha_{\mathrm{off},\mathrm{on}}=0$
  (\cref{prop:th-mv}(c)).
  \item \emph{Classifier error.} The online stream is measured with false positives
  and false negatives. Suppose that, under the null, the offline intensity depends
  only on the offline history and the calendar. Then it is predictable for every
  filtration that contains these, so classifier errors that are independent of the
  offline stream cannot \emph{create} online-to-offline causality. Under the
  alternative they dilute the signal, attenuating the estimated cross-excitation and
  the power of the test.
  \item \emph{Under-recording with reverse excitation.} Suppose offline incidents are
  recorded with $p<1$, trigger online volume ($\alpha_{\mathrm{on},\mathrm{off}}>0$)
  and self-excite ($\alpha_{\mathrm{off},\mathrm{off}}>0$). Relative to the observed
  history $\mathcal G$, the intensity of recorded offline incidents is
  $p\,\E[\lambda_{\mathrm{off}}(t)\mid\mathcal G_{t-}]$. This depends on the online
  history whenever that history is informative about \emph{unrecorded} offline
  incidents. Online volume may then predict recorded offline incidents with
  $\alpha_{\mathrm{off},\mathrm{on}}=0$. We do not formalise the size of this effect.
  \item \emph{Timestamps.} If offline incidents are dated by report time rather than
  occurrence time, online reactions to an incident can precede its record and appear
  as online-to-offline excitation.
\end{enumerate}
\end{remark}
